\documentclass[11pt]{article}

\usepackage{fullpage}
\usepackage{amsmath, amssymb, bm, cite, epsfig, psfrag}
\usepackage{graphicx}
\usepackage{float}
\usepackage{amsthm}
\usepackage{amsfonts}
\usepackage{listings}
\usepackage{cite}
\usepackage{hyperref}
\usepackage{tikz}
\usepackage{enumerate}
\usepackage{mcode}
\usetikzlibrary{shapes,arrows}
%\usetikzlibrary{dsp,chains}

%\restylefloat{figure}
%\theoremstyle{plain}      \newtheorem{theorem}{Theorem}
%\theoremstyle{definition} \newtheorem{definition}{Definition}

\def\del{\partial}
\def\ds{\displaystyle}
\def\ts{\textstyle}
\def\beq{\begin{equation}}
\def\eeq{\end{equation}}
\def\beqa{\begin{eqnarray}}
\def\eeqa{\end{eqnarray}}
\def\beqan{\begin{eqnarray*}}
\def\eeqan{\end{eqnarray*}}
\def\nn{\nonumber}
\def\binomial{\mathop{\mathrm{binomial}}}
\def\half{{\ts\frac{1}{2}}}
\def\Half{{\frac{1}{2}}}
\def\N{{\mathbb{N}}}
\def\Z{{\mathbb{Z}}}
\def\Q{{\mathbb{Q}}}
\def\R{{\mathbb{R}}}
\def\C{{\mathbb{C}}}
\def\argmin{\mathop{\mathrm{arg\,min}}}
\def\argmax{\mathop{\mathrm{arg\,max}}}
%\def\span{\mathop{\mathrm{span}}}
\def\diag{\mathop{\mathrm{diag}}}
\def\x{\times}
\def\limn{\lim_{n \rightarrow \infty}}
\def\liminfn{\liminf_{n \rightarrow \infty}}
\def\limsupn{\limsup_{n \rightarrow \infty}}
\def\GV{Guo and Verd{\'u}}
\def\MID{\,|\,}
\def\MIDD{\,;\,}

\newtheorem{proposition}{Proposition}
\newtheorem{definition}{Definition}
\newtheorem{theorem}{Theorem}
\newtheorem{lemma}{Lemma}
\newtheorem{corollary}{Corollary}
\newtheorem{assumption}{Assumption}
\newtheorem{claim}{Claim}
\def\qed{\mbox{} \hfill $\Box$}
\setlength{\unitlength}{1mm}

\def\bhat{\widehat{b}}
\def\ehat{\widehat{e}}
\def\hhat{\widehat{h}}

\def\phat{\widehat{p}}
\def\qhat{\widehat{q}}
\def\rhat{\widehat{r}}
\def\shat{\widehat{s}}
\def\uhat{\widehat{u}}
\def\ubar{\overline{u}}
\def\vhat{\widehat{v}}
\def\xhat{\widehat{x}}
\def\xbar{\overline{x}}
\def\zhat{\widehat{z}}
\def\zbar{\overline{z}}
\def\la{\leftarrow}
\def\ra{\rightarrow}
\def\MSE{\mbox{\small \sffamily MSE}}
\def\SNR{\mbox{\small \sffamily SNR}}
\def\SINR{\mbox{\small \sffamily SINR}}
\def\arr{\rightarrow}
\def\Exp{\mathbb{E}}
\def\var{\mbox{var}}
\def\Tr{\mbox{Tr}}
\def\tm1{t\! - \! 1}
\def\tp1{t\! + \! 1}

\def\Xset{{\cal X}}

\newcommand{\one}{\mathbf{1}}
\newcommand{\abf}{\mathbf{a}}
\newcommand{\bbf}{\mathbf{b}}
\newcommand{\dbf}{\mathbf{d}}
\newcommand{\ebf}{\mathbf{e}}
\newcommand{\gbf}{\mathbf{g}}
\newcommand{\hbf}{\mathbf{h}}
\newcommand{\hbfhat}{\widehat{\mathbf{h}}}
\newcommand{\pbf}{\mathbf{p}}
\newcommand{\pbfhat}{\widehat{\mathbf{p}}}
\newcommand{\qbf}{\mathbf{q}}
\newcommand{\qbfhat}{\widehat{\mathbf{q}}}
\newcommand{\rbf}{\mathbf{r}}
\newcommand{\rbfhat}{\widehat{\mathbf{r}}}
\newcommand{\sbf}{\mathbf{s}}
\newcommand{\sbfhat}{\widehat{\mathbf{s}}}
\newcommand{\ubf}{\mathbf{u}}
\newcommand{\ubfhat}{\widehat{\mathbf{u}}}
\newcommand{\utildebf}{\tilde{\mathbf{u}}}
\newcommand{\vbf}{\mathbf{v}}
\newcommand{\vbfhat}{\widehat{\mathbf{v}}}
\newcommand{\wbf}{\mathbf{w}}
\newcommand{\wbfhat}{\widehat{\mathbf{w}}}
\newcommand{\xbf}{\mathbf{x}}
\newcommand{\xbfhat}{\widehat{\mathbf{x}}}
\newcommand{\xbfbar}{\overline{\mathbf{x}}}
\newcommand{\ybf}{\mathbf{y}}
\newcommand{\zbf}{\mathbf{z}}
\newcommand{\zbfbar}{\overline{\mathbf{z}}}
\newcommand{\zbfhat}{\widehat{\mathbf{z}}}
\newcommand{\Ahat}{\widehat{A}}
\newcommand{\Abf}{\mathbf{A}}
\newcommand{\Bbf}{\mathbf{B}}
\newcommand{\Cbf}{\mathbf{C}}
\newcommand{\Bbfhat}{\widehat{\mathbf{B}}}
\newcommand{\Dbf}{\mathbf{D}}
\newcommand{\Gbf}{\mathbf{G}}
\newcommand{\Fbar}{\overline{F}}
\def\Hhat{\widehat{H}}
\newcommand{\Hbf}{\mathbf{H}}
\newcommand{\Kbf}{\mathbf{K}}
\newcommand{\Pbf}{\mathbf{P}}
\newcommand{\Phat}{\widehat{P}}
\newcommand{\Qbf}{\mathbf{Q}}
\newcommand{\Rbf}{\mathbf{R}}
\newcommand{\Rhat}{\widehat{R}}
\newcommand{\Sbf}{\mathbf{S}}
\newcommand{\Ubf}{\mathbf{U}}
\newcommand{\Vbf}{\mathbf{V}}
\newcommand{\Wbf}{\mathbf{W}}
\newcommand{\Xhat}{\widehat{X}}
\newcommand{\Xbf}{\mathbf{X}}
\newcommand{\Ybf}{\mathbf{Y}}
\newcommand{\Zbf}{\mathbf{Z}}
\newcommand{\Zhat}{\widehat{Z}}
\newcommand{\Zbfhat}{\widehat{\mathbf{Z}}}
\def\alphabf{{\boldsymbol \alpha}}
\def\betabf{{\boldsymbol \beta}}
\def\betahat{\widehat{\beta}}
\def\mubf{{\boldsymbol \mu}}
\def\lambdabf{{\boldsymbol \lambda}}
\def\etabf{{\boldsymbol \eta}}
\def\xibf{{\boldsymbol \xi}}
\def\taubf{{\boldsymbol \tau}}
\def\sigmahat{{\widehat{\sigma}}}
\def\thetabf{{\bm{\theta}}}
\def\thetabfhat{{\widehat{\bm{\theta}}}}
\def\thetahat{{\widehat{\theta}}}
\def\mubar{\overline{\mu}}
\def\muavg{\mu}
\def\sigbf{\bm{\sigma}}
\def\etal{\emph{et al.}}
\def\Ggothic{\mathfrak{G}}
\def\Pset{{\mathcal P}}
\newcommand{\bigCond}[2]{\bigl({#1} \!\bigm\vert\! {#2} \bigr)}
\newcommand{\BigCond}[2]{\Bigl({#1} \!\Bigm\vert\! {#2} \Bigr)}


\begin{document}

\title{EE6333:  Homework 3\\
Bayesian Estimation}
\author{Prof.\ Sundeep Rangan}
\date{}

\maketitle

Send all submissions to the TA,
Mustafa Riza Akdeniz {\tt akdeniza@nyu.edu}, or submit your answers
to NYU Classes.

\begin{enumerate}


\item \label{prob:pexpH}
Suppose that $p(x|y) = C(y)e^{-H(x,y)}$ 
for some functions $C(y)$ and $H(x,y)$.  
\begin{enumerate}[(a)]
\item Show that the MAP estimate of $x$ given $y$ is given by
\[
    \xhat = \argmin_x H(x,y).
\]
\item Let
\[
    Z_r := \int x^r \exp(-H(x,y))dx.
\]
Show that the MMSE estimate and its variance is given by
\[
    \Exp(x|y) = \frac{Z_1}{Z_0}, \quad \var(x|y) = \frac{Z_2}{Z_0} - \frac{Z_1^2}{Z_0^2}.
\]
\item Show that if $p(x,y) = C(y)e^{-H(x,y)}$, then the formulae
 in (a) and (b) still hold.
\end{enumerate}

\item Suppose that $\theta$ is exponentially distributed with
\[
    p(\theta) = \frac{1}{\lambda}e^{-\theta/\lambda}, 
    \quad \theta \geq 0,
\]
and we have an observation
\[
    x = \theta + w, \quad w \sim {\mathcal N}(0,\sigma^2),
\]
where $w$ is independent of $\theta$.
\begin{enumerate}[(a)]
\item What is the ML estimate of $\theta$ subject to $\theta \geq 0$?


\item Find a function $H(\theta,x)$ such that
\[
    p(\theta|x) = C(x)e^{-H(\theta,x)},
\]
for some $C(\theta)$.

\item Use the previous problem to compute the MAP estimate.
Note that to apply the constraint $\theta \geq 0$, you should
separately consider the cases $\theta > 0$ and $\theta = 0$.

\item Read out a \emph{truncated normal distribution} in
\begin{quote}
    \url{http://en.wikipedia.org/wiki/Truncated_normal_distribution}
\end{quote}
Show that the posterior density $p(\theta|x)$ is a 
truncated normal distribution.  Use the equations
on that page to derive an expression for $\Exp(\theta|x)$.
Your expression will have a $Q$-function in it.


\item Use MATLAB to create a plot with the ML, MAP and MMSE estimate
$\thetahat$ for $x \in [-2,3]$ for $\lambda=1$ and $\sigma^2=0.1$.
Plot all three estimators in the same figure and use the \mcode{legend}
command to identify the estimators.

\end{enumerate}

\item For a scalar
random variable $\theta$, let $\thetahat$ be the estimator that
minimizes the expected Bayes risk
\[
    \thetahat = \argmin_{\thetahat} \Exp\left[ L(\theta,\thetahat) \mid x \right],
\]
 for the so-called the $L_1$ loss function defined by
\[
    L(\theta,\thetahat) := |\theta-\thetahat|.
\]
Thus, $\thetahat$ minimizes the expected absolute difference.  Show that if
there exists a continuous posterior density $p(\theta|x)$, then $\thetahat$ satisfies
\[
    P(\theta > \thetahat|x) = P(\theta < \thetahat|x).
\]
That is, $\thetahat$ is the value where $\theta$ is equally likely to be above and below
$\thetahat$, which is the \emph{median} of the density.

Hint:  Write
\[
    J(\thetahat) := \Exp\left[ L(\theta,\thetahat) \mid x) \right] =
        \int_{-\infty}^\infty |\theta-\thetahat|p(\theta|x)d\theta.
\]
Then, break this integral into two parts:  one for $\theta > \thetahat$ and one for
$\theta < \thetahat$.  Then, take the derivative $J'(\thetahat)$ using Leibnitz rule:
\begin{quote}
    \url{http://en.wikipedia.org/wiki/Leibniz_integral_rule}
\end{quote}


\item A received signal has some power average power level $\theta$.  We make
$N$ measurements, $x_i$. Conditional on $\theta$, the power measurements are independent
and exponentially distributed with $\Exp(x_i|\theta) = 1/\theta$.
\begin{enumerate}[(a)]
\item What is the ML estimate of $\theta$ given the $N$ measurements $\xbf=(x_1,\ldots,x_N)$.

\item Look up the Gamma distribution.  Suppose that $\theta$ has a Gamma prior:
$\theta \sim \Gamma(\alpha,\beta)$.  Show that the posterior density of $\theta$ given $\xbf$
is also a Gamma with $\theta_{|\xbf} \sim \Gamma(\alpha',\beta')$ for some $\alpha'$ and $\beta'$.  What are $\alpha'$ and $\beta'$?  This shows that the Gamma distribution
 is the conjugate prior to the exponential.  Feel free to use any facts you can find on
wikipedia or any other source.  You can state those results without proof.

\item Using the results from above, write the posterior mean
 and variance $\Exp(\theta|\xbf)$ and $\var(\theta|\xbf)$ in
 terms of the prior mean 
 $\mu = \Exp(\theta)$ and variance $\tau = \var(\theta)$ 
 as well as $\xbf$.

\item What is the limit,
\[
    \lim_{\tau \arr \infty} \Exp(\theta|\xbf).
\]
How does this compare to the MLE in part (a)?

\end{enumerate}

\end{enumerate}

\end{document}

\int_\xhat^\infy \log(x/xhat) p(x|y)dx + \int \log(xhat/x)

P(x > xhat)/xhat = 1/xhat
